\cotefile{156-8}
\scanpage{42}
\mspage{33}

J'en viens maintenant à la question d'une fonction genre ``support''
\[
X\longmapsto |X|\qquad \mathcal M\longrightarrow\unsure{\mathscr P|\mathcal P}(P),
\]
% margin note written slanted
\margin{On \unsure{pose|veut} pour toute partie $A$ de $\mathcal M$, $|A|=\bigcup_{X\in A}|X|$.}
$P$ \unsure{est} un ensemble (dont les éléments jouent le rôle de ``points''). Voici le minimum auquel \unsure{penser} immédiatement (Comparer \unsure{CF|Cf} \uline{VII}, p.~73) :

\margin{concerne $\ll$}
\[
X\ll Y\Longrightarrow |X|\subset |Y|\tag*{(i)}
\]
\deleted{\unsure{(ii)} $X\unsure{\neq}0$}
\margin{concerne $\leq$\unsure{.} (\uline{\unsure{anti-disjonc}} \uline{\unsure{par (iii)}})}
\[
\left[\quad \forall X\in\mathcal M,\qquad |X|^{\circ}\overset{\mathrm{def}}{=}|X|\setminus|\partial X|\ne\varnothing\quad\right]\tag*{(ii)}
\]
\margin{concerne $|\circ|$}
\[
\forall X,Y\in\mathcal M,\qquad X\mathrel{|\circ|}Y\Longleftrightarrow |X|^{\circ}\cap |Y|^{\circ}=\varnothing\tag*{(iii)}
\]
\[
\begin{aligned}
&\text{(Pour ``\unsure{normaliser}'' }X\text{)}\quad |\mathcal M|=P,\text{ et les }|X|^{\circ}\text{ séparent}\\
&\text{les points de }X\unsure{,.}
\end{aligned}\tag*{(iv)}
\]

\uline{NB} La condition (i) équivaut à dire que $X\mapsto|X|$ est croissant pour $\ll$, ou encore, pour $\overset{\circ}{\ll}$ et $\leq$. Si les $|X|$ séparent les points de $\textcircled{X}$, en est-il de même des $|X|^{\circ}$ ? Soit $x\in|X|$, $y\notin|X|$. On voudrait trouver $Y\leq X$ avec $x\in|Y|^{\circ}$, bien sûr on aura $y\notin|Y|^{\circ}$ puisque $|Y|^{\circ}\subset|Y|\subset|X|$. Donc la question est de \deleted{\unsure{trouver}} savoir si $\exists Y\in\widetilde X$ avec $x\in|Y|^{\circ}$. Donc on est amené à \unsure{renforcer} (ii) en
\margin{$\unsure{P}\longrightarrow$}
\[
\begin{aligned}
&\forall X\in\mathcal M,\text{ la famille des }|Y|,\ Y\in\widetilde X,\text{ forme}\\
&\text{une figure \inserted{ens.} \deleted{de} (élémentaire) dans }P,\text{ indexée}\\
&\text{par }\widetilde X,\text{ de support }|X|.
\end{aligned}\tag*{(ii bis)}
\]
\deleted{Cela implique \unsure{gratuitement} que} \inserted{\unsure{On} sait déjà (indép. de (ii bis)) que}
\[
\widetilde X\longrightarrow\mathscr P(P),\quad Y\longmapsto|Y|
\]
est injective, car si $Y,Y'\in\widetilde X$ avec $Y\ne Y'$, on a $Y\mathrel{|\circ|}Y'$ donc par (iii) $|Y|^{\circ}\cap|Y'|^{\circ}=\varnothing$, donc par (ii) $\unsure{Y\ne Y'|\lvert Y\rvert\ne\lvert Y'\rvert}$.
\deleted{\unsure{Comme on sait Je dis que}}
% archive sheet number: 42
